\section{The reading list}
\label{sec:readinglist}

The list contains five items: four journal articles and one short monograph.
Table~\ref{tab:readinglist} gives the overview; the paragraphs below explain
what each item contributes to the arc and why I chose it over its neighbours.

\begin{table}[t]
  \centering
  \small
  \caption{The reading list. Each item marks a shift in where the reference
  behaviour for ``clean data'' comes from.}
  \label{tab:readinglist}
  \begin{tabular}{@{}p{4.6cm}lp{6.6cm}@{}}
    \toprule
    Work & Year & Place in the lineage \\
    \midrule
    \citet{grubbs1969}, \work{Technometrics} & 1969 &
      consolidation of classical outlier tests under the normal model \\
    \citet{iglewicz1993}, ASQC booklet & 1993 &
      robust reformulation for practitioners; the modified z-score and the
      3.5 rule \\
    \citet{shafer2000}, \work{J.\ Atmos.\ Oceanic Technol.} & 2000 &
      first end-to-end QA system for a large automated surface network \\
    \citet{campbell2013}, \work{BioScience} & 2013 &
      generalised QA/QC framework for streaming ecological sensor data \\
    \citet{leigh2019}, \work{Sci.\ Total Environ.} & 2019 &
      statistical anomaly detection framework, evaluated on labelled in situ
      water quality data \\
    \bottomrule
  \end{tabular}
\end{table}

\citet{grubbs1969} is the natural anchor for the classical period. The paper
is itself a work of consolidation: it gathers Grubbs' own 1950 results, the
masking analysis of \citet{pearson1936}, and two decades of tables and
significance levels into a single practical guide to testing outlying
observations in samples. I chose it over the earlier \citet{grubbs1950},
which is mathematically richer but less readable, because the 1969 paper is
the version the applied world actually cites, and part of what I want to
study is how ideas travel.

\citet{iglewicz1993} is the only non-journal item, a sixteenth volume in a
series of ASQC booklets written to put usable statistics in the hands of
quality engineers. It earns its place for one reason: it is the direct
source, cited by name in modern sensor QC documentation including that of the
system I worked on, of the modified z-score
$M_i = 0.6745\,(x_i - \tilde{x})/\mathrm{MAD}$ and the recommendation to flag
observations with $\lvert M_i\rvert > 3.5$. When a constant appears
unexplained in production code thirty years later, the honest move is to read
the book it came from.

\citet{shafer2000} documents the quality assurance system of the Oklahoma
Mesonet, the first state-wide automated surface network to take QA seriously
as an operational discipline rather than an afterthought. The paper fixed the
canonical test battery, range, step, persistence, like-instrument and spatial
tests, and, just as important, the surrounding architecture of calibration,
rotation, flag taxonomy and human review. Nearly every network QA paper since
either cites it or silently reimplements it.

\citet{campbell2013} lifted that operational knowledge out of meteorology and
into ecology at the moment cheap sensors were flooding ecological science
with data. It is the framework paper of the streaming era: it separates
quality assurance from quality control, catalogues sensor fault types,
formalises the pipeline from field to archive, and insists on the
human-in-the-loop and on versioned, flagged archives rather than silent
deletion. I chose it over software-specific descendants because it states the
design principles those tools implement.

\citet{leigh2019} closes the list at the research frontier that my thesis
topic touches. Written by a team spanning a state environmental agency and a
statistics centre of excellence, it defines anomaly types for high-frequency
water quality series, compares regression-based, feature-based and rule-based
detectors on labelled data from river sensors, and reports honest,
type-resolved performance numbers. Among candidate frontier papers I
preferred it to a pure machine learning survey because it keeps the
operational framing, the flags still have to mean something to a hydrologist,
while importing serious statistical machinery.

Around this spine I cite supporting works where the roles demand them:
\citet{peirce1852} and \citet{chauvenet1863} for the astronomical prehistory,
\citet{anscombe1960} and \citet{hampel1974} for the conceptual turn to
insurance and robustness, \citet{rosner1983} for the repair of sequential
testing, \citet{gandin1988} for the parallel tradition in numerical weather
prediction, \citet{fiebrich2010} and \citet{estevez2011} for the
consolidation of network QA practice, \citet{taylor2013} for observatory
scale automation, \citet{talagala2019} for the feature-based line,
\citet{jones2022} and \citet{schmidt2023} for the software generation, and
\citet{blazquez2021} for the machine learning panorama. Two exclusions were
deliberate. I left out the WMO guidance documents because they codify
practice rather than originate it, and I left out deep learning anomaly
detection papers because, as Section~\ref{sec:futurists} argues, their
contact with operational environmental QC is still thin enough that reading
them here would be speculation rather than lineage.
